\metatitle{Thom polynomials}
\metadescription{An introduction to Thom polynomials, their expression in Chern classes, and their Schur function expansions.}
\metakeywords{Thom polynomials,Schur functions,supersymmetric Schur functions,Chern classes,singularity theory}

\section[thomPolynomials]{Thom polynomials}

\begin{polydata}{thomPolynomial}
  Name   & Thom polynomials \\
  Space  & SuperSym \\
  Basis  & False \\
  Rating & 1 \\
  Bib    & Thom1956 \\
  Year   & 1956 \\
  Keywords & schur-positive,chern-classes,singularity \\
  Symbol & $\mathcal{T}^{\eta}$ \\
  Category & Other \\
  PositiveIn & hookSchur | PragaczWeber2007 | scope=stable singularities with positive sum \\
\end{polydata}

The \defin{Thom polynomial} of a singularity class $\eta$ is a universal
polynomial in Chern classes whose value gives the cohomology class of the
$\eta$-locus of a sufficiently generic map.  The subject goes back to
\name{René Thom}'s work on singularities of differentiable maps
\cite{Thom1956}.

For maps $f : M \to N$ of complex manifolds, the relevant Chern classes are
usually the Chern classes of the virtual bundle $f^\ast TN - TM$.  If we pass
to Chern roots, this means that Thom polynomials are naturally symmetric
polynomials in a \hyperref[hookSchur]{difference of alphabets}.  This is why
\hyperref[hookSchur]{supersymmetric Schur functions}, also called Schur
functions in a difference of alphabets, are a natural basis for them.

\name{Piotr Pragacz} used Schur and supersymmetric Schur functions to compute
and organize Thom polynomials for several singularity classes
\cite{Pragacz2005x,Pragacz2007,Pragacz2008x}.  One advantage of the Schur
basis is that it can reveal structural restrictions and recurrences in the
partitions that occur, which are much harder to see from monomials in Chern
classes.

\name{Piotr Pragacz} and \name{Andrzej Weber} proved positivity results for
Schur function expansions of Thom polynomials \cite{PragaczWeber2007}.  In
particular, for stable singularities satisfying their positivity hypotheses,
the Schur expansion has non-negative coefficients.  This places Thom
polynomials in the same general circle of geometric Schur-positivity
phenomena as degeneracy-locus formulas and Schubert calculus.

